\section{Тестирование}

\subsection{Модульное тестирование}

\headingtextsep

Модульное тестирование предназначено для проверки корректности изолированных частей серверной и клиентской логики. Основной задачей этого этапа является подтверждение того, что отдельные функции, методы и небольшие прикладные компоненты работают в соответствии с заданными правилами обработки данных.

На серверной стороне модульными тестами должны покрываться операции создания и изменения задач, вычисление зависимостей и признаков блокировки, проверка прав доступа, обработка агентных разрешений, формирование записей аудита и правила работы со связями между сущностями. Отдельно должны проверяться преобразования данных, валидация входных значений и ограничения, связанные со статусами, сроками и назначением исполнителей.

На клиентской стороне модульные тесты должны применяться для проверки логики отображения данных, переходов между состояниями интерфейса, валидации форм, обработки ошибок и преобразования ответов API в пользовательские модели представления. Особое внимание должно уделяться экранам задач, документации, заметок, администрирования и поисковой выдачи.

Результатом модульного тестирования должно быть подтверждение того, что отдельные программные элементы корректно работают на допустимых и ошибочных наборах входных данных и не нарушают базовые инварианты предметной области.

\textheadingsep
\subsection{Интеграционное тестирование}

\headingtextsep

Интеграционное тестирование предназначено для проверки совместной работы нескольких подсистем и корректности обмена данными между ними. На этом уровне должны проверяться сценарии взаимодействия веб-интерфейса с API, API с базой данных, API с файловым хранилищем, а также работа агентного доступа и внешних интеграций.

Интеграционные тесты должны охватывать создание и изменение задач с сохранением данных в PostgreSQL, загрузку файловых ресурсов с вынесением бинарного содержимого во внешнее хранилище, чтение и обновление документации и заметок, выпуск и отзыв агентных ключей, а также фиксацию действий в журнале аудита. Для поиска должны проверяться операции индексации, полнотекстового поиска и обработки семантических запросов по тем данным, которые были ранее загружены в систему.

Отдельная группа интеграционных тестов должна проверять работу механизмов аутентификации и авторизации: вход пользователя, применение ролей, разграничение доступа к административным функциям, ограничение области действия агентных ключей и корректную обработку отказов в доступе.

При наличии подключённой внешней интеграции должны выполняться проверки связывания внутренних сущностей с объектами GitHub, сохранения параметров интеграции и корректного отображения связанных внешних данных внутри системы.

\textheadingsep
\subsection{Сценарное тестирование}

\headingtextsep

Сценарное тестирование используется для проверки ключевых пользовательских и административных потоков работы системы в условиях, близких к реальной эксплуатации. Оно должно подтверждать, что пользователь может выполнить законченную рабочую задачу без нарушения логики интерфейса и бизнес-правил.

В рамках сценарного тестирования должны проверяться следующие типовые сценарии: вход пользователя в систему, выбор рабочего пространства и проекта, создание группы задач и связанных задач, назначение исполнителя, изменение статуса задачи, добавление заметки, создание страницы документации, загрузка связанного файла, выполнение поиска по задачам и документам, просмотр журнала активности, а также выполнение базовых административных операций по управлению участниками и агентами.

Отдельно должны проверяться сценарии машинного взаимодействия: подключение агента по выданному ключу, получение разрешённых данных по проекту, создание заметки агентом, ограниченное обновление статуса задачи и фиксация результатов агентной сессии в журнале системы.

Для сценариев с ошибками должны выполняться проверки некорректного ввода, отказа в доступе, работы с просроченным или отозванным агентным ключом, попыток обращения к недоступным объектам и обработки временной недоступности внешних сервисов.

\textheadingsep
\subsection{Критерии приемки}

\headingtextsep

Система может считаться готовой к приемке при выполнении набора функциональных и качественных критериев. Все обязательные сценарии работы с задачами, документацией, заметками, файлами, поиском, агентным доступом и административными функциями должны выполняться без критических ошибок и без потери данных.

Приемка должна подтверждать успешное прохождение модульных, интеграционных и сценарных тестов по основным функциям системы. Для каждой проверяемой операции должно быть установлено, что данные сохраняются корректно, связи между сущностями не нарушаются, права доступа соблюдаются, а журнал аудита фиксирует значимые действия.

Для файловых ресурсов должно быть подтверждено корректное сохранение метаданных в базе данных, размещение бинарного содержимого во внешнем хранилище и возможность последующего чтения или скачивания файла в соответствии с правами доступа. Для поисковых функций должно быть подтверждено, что система находит релевантные задачи, документы, заметки и файловые ресурсы по их содержимому, имени или связанным признакам.

Дополнительным критерием приемки является успешная работа системы после повторного запуска сервисов и корректное применение конфигурации локальной среды разработки. При наличии подключённых внешних интеграций должно быть подтверждено, что система корректно работает с их данными и не теряет внутреннюю согласованность при временных ошибках внешней стороны.